\documentclass[11pt]{article}

\usepackage[a4paper,margin=1in]{geometry}
\usepackage[T1]{fontenc}
\usepackage[utf8]{inputenc}
\usepackage{amsmath,amssymb,amsthm}
\usepackage{mathtools}
\usepackage{enumitem}
\usepackage{booktabs}
\usepackage{array}
\usepackage{tabularx}
\usepackage{xcolor}
\usepackage{hyperref}
\usepackage{cleveref}

\hypersetup{colorlinks=true,linkcolor=blue!50!black,citecolor=blue!50!black,urlcolor=blue!50!black}

\newtheorem*{defthm}{Definition}
\newtheorem*{thmthm}{Theorem}
\newtheorem{proposition}{Proposition}
\crefname{proposition}{Proposition}{Propositions}
\newtheorem*{remthm}{Remark}

\newcommand{\CP}{\mathbb{CP}}
\newcommand{\type}{\mathrm{type}}
\newcommand{\level}{\mathrm{level}}
\newcommand{\verify}{\textcolor{red!70!black}{\textsf{[verify]}}}
\newcolumntype{L}{>{\raggedright\arraybackslash}X}

\title{\bf Maubach Subdivision on $\CP^2$\\[4pt]\large Related Work}
\author{Vinicius Mello}
\date{\today}

\begin{document}
\maketitle

\begin{abstract}
This document surveys the literature related to a planned article built on
the technical report \emph{Maubach Subdivision on $\CP^2$}. The article has
three strands: (i) the combinatorial description of Maubach's subdivision as
a \emph{stellar subdivision scheme}~\cite{mello2006} and its adaptation to
arbitrary compatible initial meshes; (ii) pointerless LPT
codes~\cite{atalaymount2007} and their extension to non-cubical seeds; and
(iii) the triangulation of Riemann surfaces as codimension-2 loci of $\CP^2$
as an application. For each strand we list the closest prior work and state
precisely how our approach compares. Along the way we confirm
computationally that Gaifullin's 15-vertex triangulation of
$\CP^2$~\cite{gaifullin2009} is \emph{balanced}, i.e.\ properly
$5$-colorable, as Gaifullin himself states (with minimality among balanced
triangulations of $\CP^2$, \cite[\S2]{gaifullin2009}). Its Maubach
compatibility is therefore an instance of the classical
``coloring $\Rightarrow$ reflected'' mechanism. Items marked \verify{} were not checked against the original source.
\end{abstract}

\tableofcontents

% =====================================================================
\section{Summary and positioning}
% =====================================================================

\begin{enumerate}[leftmargin=*,itemsep=4pt]
\item \textbf{The direct competitor is Diening, Gehring and Storn
  (DGS)~\cite{dgs2025}.} They initialize Maubach's routine on \emph{any}
  conforming triangulation in any dimension by a proper $(N{+}1)$-coloring
  of the vertices, $N\ge n$. The coloring is computed greedily in linear
  time, with $N$ bounded by the maximal vertex degree. For $N=n$ the coloring
  is exactly the ``global flag/rank'' mechanism behind Theorems 9 and 10 of
  the thesis. For $N>n$ each $n$-simplex is treated as a face of a virtual
  $N$-simplex. The article must position itself explicitly with respect to
  DGS.

\item \textbf{Gaifullin's $\CP^2$ is balanced}
  (\Cref{prop:gaifullin}). The Maubach ordering found by the search of the
  report is precisely a proper $5$-coloring of the vertices. This contradicts
  the report's statement that the mesh has no global rank. What is true is
  that the rank is not given \emph{by construction}. This yields a clean
  statement for the article: balanced $\Rightarrow$ Maubach-compatible with
  $i=j$ on every shared facet~\cite{traxler1997,dgs2025}.

\item \textbf{Kühnel's $\CP^2_9$ cannot be balanced.} It is
  $3$-neighborly, so its $1$-skeleton is $K_9$ and needs $9$ colors. In DGS
  terms it is the case $N=8>n=4$.

\item \textbf{Theorems 9 and 10 of the thesis are not new as facts.}
  Products of simplices and barycentric subdivisions are the textbook
  balanced complexes, and ``balanced $\Rightarrow$ reflected'' is already in
  Traxler~\cite{traxler1997}. The pre-refinement of
  Kossaczký~\cite{kossaczky1994} and Stevenson~\cite[App.~A]{stevenson2008}
  is essentially the barycentric subdivision without the edge-midpoint
  level, with $(n{+}1)!/2$ cells instead of $(n{+}1)!$. What belongs to the
  thesis is the \emph{formalism}:
  \begin{itemize}[itemsep=0pt]
  \item a scheme is an invariant plus an algorithm;
  \item the invariant is proved inductively, level by level (Lemmas 5--7,
    Theorem 8);
  \item Mitchell's facet-based and Maubach's edge-based schemes are unified
    in one stellar language.
  \end{itemize}

\item \textbf{The LPT strand has strong analogues but no equivalent.}
  Existing work includes t8code/TM-index~\cite{bursteddeholke2016},
  concurrent binary trees~\cite{dupuy2020}, diamond
  hierarchies~\cite{weissdefloriani2011} and permutahedral
  encodings~\cite{bkw2023}. None of them combines conforming bisection,
  arbitrary dimension, a non-cubical seed and no pointers.

\item \textbf{The application occupies a real gap.} We triangulate the curve
  \emph{intrinsically in $\CP^2$}: compact, with no charts, no branch cuts
  and no special treatment of points at infinity. We do this by
  combinatorial continuation over a Maubach-refined mesh of $\CP^2$ itself.
  The weakness to state openly is that the output is \emph{not}
  topologically certified.
\end{enumerate}

% =====================================================================
\section{Bisection refinement: newest vertex and Maubach}
\label{sec:bisection}
% =====================================================================

\subsection{The classical lineage}

\begin{center}
\small
\begin{tabularx}{\textwidth}{@{}>{\raggedright\arraybackslash}p{3.6cm}LL@{}}
\toprule
Work & Contribution & Requirement on the initial mesh \\
\midrule
Rivara~\cite{rivara1984,rivara1991} & Longest-edge bisection (2D/3D/$n$D) & None combinatorial; the choice is geometric \\
Mitchell~\cite{mitchell1988,mitchell1991} & Newest vertex bisection (NVB), 2D & Compatible refinement edges; always exist in 2D~\cite{bdd2004} \\
Bänsch~\cite{bansch1991}, Kossaczký~\cite{kossaczky1994} & NVB in 3D & Kossaczký pre-refines the initial mesh \\
\textbf{Maubach}~\cite{maubach1995} & $n$D NVB: bisect $[v_0,v_\gamma]$, cyclic tag & Grids ``generated by reflection'' \\
\textbf{Traxler}~\cite{traxler1997} & Equivalent formulation & Reflected partition. For simply connected domains: exists iff every interior $(n{-}2)$-face lies in an even number of simplices \\
Arnold--Mukherjee--Pouly~\cite{amp2000} & Marked tetrahedral bisection (5 types) & Any 3D mesh, but only uniform refinements of level $\equiv 0 \pmod n$ are conforming \\
\textbf{Stevenson}~\cite{stevenson2008} & Closure estimate in $n$D & Matching-neighbour condition (b), weaker than Traxler's. Open whether every 3D mesh satisfies it. App.~A: pre-refinement into $\tfrac12(n{+}1)!$ simplices always does \\
Alkämper--Gaspoz--Klöfkorn~\cite{agk2018} & Weak compatibility: $O(n)$ renumbering of any mesh. Termination, not optimal closure & Any \\
Belda-Ferrín et al.~\cite{belda2018,belda2022,belda2023} & $n$ marked bisection stages, then ``cast'' to reflected simplices and switch to Maubach & Any conforming mesh \\
\textbf{Diening--Gehring--Storn}~\cite{dgs2025} & Initialization by $(N{+}1)$-coloring; greedy, $O(\#\text{vertices})$; shape regularity and BDV closure & Any. $N=n$ is Traxler's condition \\
Diening--Storn--Tscherpel~\cite{dst2023} & Optimal grading; vertex generation from colors & $(n{+}1)$-colored \\
Gehring~\cite{gehring2023} & Constant in the BDV--Stevenson theorem & -- \\
Li~\cite{li2026} & First unconditional closure estimate for AMP/Bänsch & Any (3D) \\
\bottomrule
\end{tabularx}
\end{center}

Belda-Ferrín et al.\ state that, for $n>3$, ``there is not any known method to
extract a reflection structure from an arbitrary unstructured conformal
mesh''~\cite{belda2023}. DGS~\cite{dgs2025} closes this gap by coloring,
without extracting a reflection structure. Their Remark~9 notes that $(n{+}1)$-colorings do not always
exist. Their Section~3.2 therefore allows $N>n$ colors: each simplex is
embedded as a face of a virtual $N$-simplex, and the bisection vertex
generations are propagated accordingly. Their Theorem~10 shows that the
resulting routine coincides with Maubach's original one under the ordering
of their Definition~2.

\subsection{Where the thesis fits}

\paragraph{Definition 9 versus Stevenson's condition (b).}
(MaubachA) allows two same-level neighbours to have either $i=j$ or
$\{i,j\}=\{0,\type\}$. In Maubach's convention (bisection edge $[0,\type]$)
this is the analogue of Stevenson's use of
\[
T_R=(x_n,x_1,\dots,x_\gamma,x_{n-1},\dots,x_{\gamma+1},x_0)_\gamma ,
\]
the tagged simplex with the same children as $T$, in Traxler's convention
(bisection edge $x_0x_n$). There, $T$ and $T'$ are \emph{reflected
neighbours} if the vertex sequence of $T$ or of $T_R$ agrees with that of
$T'$ in all but one position~\cite[\S4]{stevenson2008}.
\verify{} the exact equivalence after translating conventions.

The structural difference is this. The thesis states an invariant that holds
at \emph{every} level and proves that the algorithm preserves it (Lemma 7,
Theorem 8). Stevenson and Traxler impose their condition on the initial mesh
only, and then prove that all uniform refinements are conforming (Stevenson's
Theorem 4.3). Both routes prove essentially the same theorem; the thesis's
proof is local, inductive and purely combinatorial. Stevenson's Remark~4.4
shows that condition (b) is also \emph{necessary}. Definition 9 can thus be
presented as a hereditary form of the weakest possible level-0 condition.

\paragraph{Theorems 9 and 10.}
Both constructions are balanced complexes in Stanley's sense. For a product
of simplices, the lattice rank $|L|$ is an $(n{+}1)$-coloring; for a
barycentric subdivision, the dimension of the flag element is. ``Balanced
$\Rightarrow$ reflected with $i=j$'' is Traxler's condition, restated as a
coloring in DGS~\S3.1. The thesis (2006) predates Stevenson (2008) but not
Traxler (1997), so priority on the general fact cannot be claimed.
Theorems 9 and 10 are best presented as instances with short proofs in the
stellar formalism. Stevenson's pre-refinement costs half as many cells as
the barycentric one, and DGS show that neither is needed. For non-balanced
meshes the fair comparison is therefore DGS with $N>n$.

\paragraph{Where there is room for a genuine contribution.}
\begin{enumerate}[itemsep=2pt]
\item \emph{Closed manifolds as seeds.} All work cited above concerns domains
  in $\mathbb R^n$, for finite elements. Triangulations of closed manifolds
  ($\CP^2$, $S^2\times S^2$) and the question of which minimal
  triangulations are balanced belong to combinatorial
  topology~\cite{joswig2002,ikn2017,venturello2019}. No bridge to NVB seems
  to exist.
\item \emph{Gaifullin's triangulation is balanced} (\Cref{prop:gaifullin}).
  \emph{Correction:} this is stated by Gaifullin himself, who also proves
  that his triangulation is vertex-minimal among balanced
  (``regularly vertex-coloured'') triangulations of
  $\CP^2$~\cite[\S2 and Prop.~2.3]{gaifullin2009}. The novelty is only the
  connection with Maubach bisection.
\item \emph{Stellar subdivision schemes as a unifying language.} Mitchell
  (facet) and Maubach (edge) fit the same form. This connects to Alexander's
  and Newman's stellar moves~\cite{lickorish1999}, to Velho's stellar
  grammars~\cite{velho2003} and $4$--$8$ meshes~\cite{velhozorin2001}
  (which are NVB in 2D), and to balanced cross-flips~\cite{ikn2017}, the
  balanced analogue of Pachner moves.
\end{enumerate}

\subsection{Colorability criteria}
\label{sec:colorability}

Joswig~\cite{joswig2002} shows that a locally strongly connected simplicial
complex is balanced if and only if its group of projectivities is trivial.
For simply connected manifolds this reduces to Traxler's criterion: every
$(n{-}2)$-face has even degree. Two elementary facts suffice for the $\CP^2$
seeds:
\begin{itemize}[itemsep=1pt]
\item a $k$-neighborly complex with $k\ge 2$ and more than $n+1$ vertices is
  never balanced, since adjacent vertices need distinct colors;
\item $\CP^2$ is simply connected, so it suffices to check that every
  triangle lies in an even number of $4$-simplices.
\end{itemize}

\begin{proposition}
\label[proposition]{prop:gaifullin}
Gaifullin's $15$-vertex, $108$-cell triangulation of $\CP^2$ is balanced.
Explicitly, its vertices split into five classes of three mutually
non-adjacent vertices. The complement of the $1$-skeleton is a disjoint
union of five triangles, and every vertex has degree $12$. In the ordering
used by \texttt{riemann\_cp2}, the position of each vertex is the same in
every cell containing it. Hence all $270$ interior facets satisfy
$\partial_i\sigma=\partial_j\omega\Rightarrow i=j$. Moreover, each of the
$240$ triangles lies in $4$ ($180$ triangles) or $6$ ($60$ triangles)
cells, and the dual graph is bipartite. Kühnel's $\CP^2_9$ is not balanced.
\end{proposition}

\begin{proof}[Verification]
Direct computation on the facet list \texttt{gaifullin\_cells\_raw} in
\nolinkurl{examples/top/riemann_cp2/riemann_cp2.cpp}:
\begin{itemize}[itemsep=1pt]
\item exhaustive backtracking finds a proper 5-coloring;
\item the per-vertex set of positions is a singleton for all 15 vertices;
\item the facet-index pairs are $(k,k)$, $54$ times for each
  $k=0,\dots,4$;
\item triangle degrees were counted directly.
\end{itemize}
For Kühnel's $\CP^2_9$, $3$-neighborliness gives a $K_9$ skeleton.
\end{proof}

\begin{remthm}
Gaifullin also shows that the automorphism group $S_4\times S_3$ acts by
Fubini--Study isometries~\cite{gaifullin2009}. This explains the measured
max/min Fubini--Study edge ratio of exactly $1.4636$ with zero variance
across all 108 cells. The same paper gives a 33-vertex subdivision on which
the moment map is simplicial, a candidate seed worth testing.
\end{remthm}

% =====================================================================
\section{Pointerless representations and neighbour finding}
\label{sec:lpt}
% =====================================================================

\begin{center}
\small
\begin{tabularx}{\textwidth}{@{}>{\raggedright\arraybackslash}p{3.4cm}L>{\raggedright\arraybackslash}p{2.1cm}p{1.3cm}p{0.9cm}@{}}
\toprule
Work & Hierarchy / idea & Seed & Conf. & Dim. \\
\midrule
Hebert~\cite{hebert1994} & Symbolic tetrahedral bisection & Cube & yes & 3 \\
Maubach~\cite{maubach1996} & Neighbour location in $n$D Maubach grids & Kuhn & yes & $n$ \\
Evans--Kirkpatrick--Townsend~\cite{ekt2001} & Right-triangulated irregular networks: binary label $\Rightarrow$ geometry & Square & yes & 2 \\
Lee--De Floriani--Samet~\cite{ldfs2001} & Constant-time neighbours for bisection tetrahedra & Cube (6 tets) & yes & 3 \\
\textbf{Atalay--Mount}~\cite{atalaymount2007} & \textbf{LPT codes}: geometry and neighbours by table lookup & Kuhn cube ($\mathrm{Sym}(n)$ orbit) & yes & $n$ \\
Weiss--De Floriani~\cite{weissdefloriani2011} & Diamond hierarchies (diamond = star of the bisection edge) & Kuhn & yes & $n$ \\
\textbf{Burstedde--Holke}~\cite{bursteddeholke2016}, t8code & TM-index (Morton SFC for Bey red refinement); \emph{forest} over an arbitrary coarse mesh & \textbf{Any} & no & 2, 3 \\
Burstedde--Wilcox--Ghattas~\cite{p4est2011} & Forest of octrees & Any (hex) & no & 2, 3 \\
Dupuy~\cite{dupuy2020} & Concurrent binary trees: implicit heap of bits; longest-edge bisection on GPU & 2D & yes & 2 \\
Boissonnat--Kachanovich--Wintraecken~\cite{bkw2023} & Permutahedral encoding of Freudenthal--Kuhn/Coxeter simplices & $\mathbb R^d$ (uniform) & yes & $d$ \\
\bottomrule
\end{tabularx}
\end{center}

\paragraph{Positioning of \texttt{glpt}.}
\begin{itemize}[leftmargin=*,itemsep=3pt]
\item \emph{Against Atalay--Mount.} Their same-orthant neighbour formula
  relies on the seed being a single $\mathrm{Sym}(n)$ orbit. A \emph{balanced}
  seed suffices. Ordering the vertices of each seed cell by color identifies
  it affinely with the Kuhn reference simplex. The LPT code inside a seed is
  then unchanged, and crossing a seed boundary preserves the local facet
  index ($i=j$). Such crossings are resolved by a seed adjacency table,
  $108\times5$ for Gaifullin. This explains \emph{why} the adaptation is
  easy. It is less easy in two cases:
  \begin{itemize}[itemsep=0pt]
  \item seeds compatible only through the $\{0,\type\}$ exception would need
    a permutation correction at crossings;
  \item DGS seeds with $N>n$ would need codes of a virtual $N$-simplex, and
    DGS's Remark~12 warns that restricting refinements of the virtual
    $N$-simplex does \emph{not} reproduce their refinement.
  \end{itemize}
  Both are future work.
\item \emph{Against t8code.} The architecture is the same (tree id
  $\approx$ \texttt{seed\_index}, per-tree linear code, inter-tree
  connectivity table). Ours uses \emph{conforming}, graded bisection
  without hanging nodes, in arbitrary dimension (4 here).
\item \emph{Against BKW.} They encode a uniform triangulation; \texttt{glpt}
  encodes an adaptive hierarchy.
\item \emph{Limitations to state.} The \texttt{DEEP\_CROSSING} case
  (neighbour refined deeper than the query) is not implemented. Storage is
  $O(\#\text{leaves})$ in a hash table.
\end{itemize}

% =====================================================================
\section{Application: triangulating Riemann surfaces}
\label{sec:riemann}
% =====================================================================

\subsection{Numerical computation of Riemann surfaces}
Deconinck and van Hoeij's Maple package \texttt{algcurves}~\cite{deconinckvanhoeij2001}
computes monodromy and period matrices of plane algebraic curves. Other
tools in the same vein include Frauendiener--Klein's Matlab codes, Sage's
\texttt{RiemannSurface} (Bruin--Sijsling--Zotine), and Kranich's certified
homotopy continuation of plane curves~\cite{kranich2016}. See
Bobenko--Klein~\cite{bobenkoklein2013} for an overview. These methods treat
the surface as a branched cover of $\CP^1$ and compute analytic invariants.
We construct a \emph{simplicial complex} of the curve in $\CP^2$, which is
complementary: our mesh can, for instance, confirm the genus through
$\chi=2-2g$.

\subsection{Visualization}
Hanson~\cite{hanson1994} renders Fermat curves in $\CP^2$ projected to
$\mathbb R^3$, using parametrizations and symmetry.
Nieser--Poelke--Polthier~\cite{nieser2010} build meshes as multivalued
graphs over $\mathbb C$ from branch points, with cuts. Kranich renders
domain-coloured Riemann surfaces on the GPU~\cite{kranich2016}. All of these
work through a projection to $\mathbb C$, with charts or cuts. Our
construction is intrinsic in $\CP^2$; the price is the chart-consistency
problem at visualization time.

\subsection{Implicit manifolds of codimension greater than one}
\begin{itemize}[leftmargin=*,itemsep=3pt]
\item \textbf{PL continuation.} Allgower--Schmidt~\cite{allgowerschmidt1985}
  approximate implicit manifolds of arbitrary codimension over
  Freudenthal--Kuhn triangulations. The textbook treatment of simplicial
  continuation is Allgower--Georg~\cite{allgowergeorg2003}. This is the
  historical lineage of our ``continuation + Kuhn''.
\item \textbf{Weigle--Banks}~\cite{weiglebanks1996}, complex-valued contour
  meshing, treat the zero set of $f:\mathbb C^2\to\mathbb C$ as a surface in
  $\mathbb R^4$, extracted simplicially on a grid. This is the most direct
  predecessor of our codimension-2 extraction, but it works on a uniform
  affine grid of $\mathbb R^4$, not on $\CP^2$.
\item Bhaniramka--Wenger--Crawfis~\cite{bwc2004} construct isosurfaces in
  any dimension, in codimension 1.
\item Castelo et al.~\cite{castelo2006}, the J1a triangulation: adaptive,
  any dimension, for implicit manifolds. This is the Brazilian lineage
  of~\cite{gomestavares1989} and a natural competitor on $n$D adaptivity.
\item \textbf{Boissonnat--Kachanovich--Wintraecken}~\cite{bkw2023} trace
  isomanifolds from a seed point by cell-to-cell continuation over
  Freudenthal--Kuhn/Coxeter triangulations, polynomially in $d$ and with
  guarantees; see also~\cite{boissonnatwintraecken2022} for topological
  correctness. This is a mandatory reference: same strategy, uniform
  triangulation, certified. Ours is adaptive and lives in $\CP^2$, but is
  not certified.
\item Ju et al.~\cite{ju2024} build adaptive simplicial grids for implicit
  complexes, including curve networks (codimension 2 in $\mathbb R^3$),
  starting from a 6-tetrahedron cube. \verify{} their exact bisection rule
  and refinement criterion.
\item Gerstner--Rumpf~\cite{gerstnerrumpf2000} and
  Pascucci~\cite{pascucci2002} extract isosurfaces on bisection
  hierarchies. The combination ``Maubach hierarchy + extraction'' is thus
  established in visualization, but only in codimension 1 and in
  $\mathbb R^3$.
\end{itemize}

\subsection{Certification}
Plantinga--Vegter~\cite{plantingavegter2004} guarantee isotopy via interval
arithmetic on quadtrees/octrees. Mourrain--Pavone~\cite{mourrainpavone2009}
subdivide in the Bernstein basis. Reuter et al.~\cite{reuter2008} use the
\emph{barycentric} Bernstein basis on simplices, the exact counterpart of
our \texttt{--bernstein} enclosure on 2-faces. Bertini\_real~\cite{bertinireal2017}
computes numerical cell decompositions of real curves and surfaces; a
complex curve is a real surface in $\mathbb R^4$ (\verify{} whether it
handles this case). Our method is heuristic and adaptive.
\texttt{--tangency-threshold} and \texttt{--repair-rounds} reduce the
bad-cell rate (best: $0.180\%$) but do not eliminate it. It should be
presented as a practical triangulation, checkable a posteriori via
$\chi=2-2g$, not as a certified one.

% =====================================================================
\section{Triangulations of $\CP^2$}
% =====================================================================

\begin{itemize}[leftmargin=*,itemsep=3pt]
\item Kühnel--Banchoff~\cite{kuhnelbanchoff1983}: $\CP^2_9$, the unique
  9-vertex triangulation. It is $3$-neighborly, hence not balanced.
\item \textbf{Gaifullin}~\cite{gaifullin2009}: 15 vertices, automorphism
  group $S_4\times S_3$ acting by Fubini--Study isometries; minimal among
  triangulations admitting a chess colouring. Balanced by
  \Cref{prop:gaifullin}. Note: the technical report cites arXiv:0904.4222
  with an incorrect title; the correct title is the one given in the
  bibliography.
\item Schwartz~\cite{schwartz2023}: trisection of $\CP^2_9$.
\item Tools used for verification: GAP with simpcomp~\cite{simpcomp}.
\end{itemize}

% =====================================================================
\section{Suggestions for the article}
% =====================================================================

\begin{enumerate}[leftmargin=*,itemsep=3pt]
\item \textbf{Central statement.} ``Maubach bisection works verbatim on any
  balanced initial triangulation, including triangulations of closed
  manifolds; pointerless LPT codes extend to such seeds by adding a seed
  index and a seed adjacency table.'' The first half is known~\cite{traxler1997,dgs2025}.
  The contributions are:
  \begin{itemize}[itemsep=0pt]
  \item the star-scheme formalism with a hereditary invariant;
  \item closed-manifold seeds (Gaifullin's balanced $\CP^2$);
  \item the LPT extension;
  \item the application.
  \end{itemize}
\item \textbf{Comparison experiments.}
  \begin{itemize}[itemsep=0pt]
  \item Kühnel with DGS initialization ($N=8$) against Gaifullin ($N=4$):
    cells, Fubini--Study quality, bad-cell rate. This would be the first use
    of DGS outside $\mathbb R^n$.
  \item Barycentric subdivision of Kühnel against Stevenson's
    pre-refinement of Kühnel against Gaifullin: cell cost and quality.
  \item A Weigle--Banks style extraction on an affine chart grid against
    intrinsic $\CP^2$ on the same curve: cells visited and artifacts near
    infinity.
  \end{itemize}
\item \textbf{Corrections to the report} before reusing text:
  \begin{itemize}[itemsep=0pt]
  \item Gaifullin \emph{has} a global rank structure (a coloring);
  \item fix the title of the Gaifullin reference;
  \item reinterpret the report's ``coloring problem on the dual graph'' as a
    search for a proper $(n{+}1)$-coloring of the vertices.
  \end{itemize}
\item \textbf{Open questions.} What is the minimal number of vertices of a
  balanced triangulation of $\CP^2$? Does Gaifullin's 33-vertex subdivision
  (simplicial moment map) give a better seed?
\end{enumerate}

\bibliographystyle{plain}
\bibliography{related_work}

\end{document}
